\documentclass{article}
 \usepackage[T1]{fontenc}
 \usepackage{array}
\usepackage{booktabs}
 \usepackage{siunitx}
\usepackage{tabularx}
\usepackage{longtable}
 \usepackage{threeparttable}
\usepackage{ragged2e}

%\setlength\tabcolsep{1cm}
 \begin{document}

\section*{5 Tables}

\subsection*{ 5.1 The array package}

 \begin{tabular}{lll}
 Animal & Food & Size \\
 dog & meat & medium \\
 horse & hay & large \\
 frog & flies & small \\
 \end{tabular}

\vspace{2em}

 \begin{tabular}{cl}
 Animal & Description \\
 dog  & The dog is a member of the genus Canis, which forms part of the  wolf-like canids, and is the most widely abundant
 terrestrial carnivore. \\
 
 cat & The cat is a domestic species of small carnivorous mammal. It is the  only domesticated species in the family Felidae and is often referred  to as the domestic cat to distinguish it from the wild members of the family. \\

 \end{tabular}

\vspace{2em}

 \begin{tabular}{cp{9cm}}
 Animal & Description \\
 dog  & The dog is a member of the genus Canis, which forms part of the  wolf-like canids, and is the most widely abundant
 terrestrial carnivore. \\
 
 cat & The cat is a domestic species of small carnivorous mammal. It is the  only domesticated species in the family Felidae and is often referred  to as the domestic cat to distinguish it from the wild members of the family. \\

 \end{tabular}

\vspace{2em}

 \begin{tabular}{*{3}{l}}
 Animal & Food & Size \\
 dog & meat & medium \\
 horse & hay & large \\
 frog & flies & small \\
 \end{tabular}

\subsection*{ 5.2 Addingrules(lines)}

 \begin{tabular}{lll}
 \toprule
 Animal & Food & Size \\
 \midrule
 dog & meat & medium \\
 horse & hay & large \\
 frog & flies & small \\
 \bottomrule
 \end{tabular}

\vspace{2em}

 \begin{tabular}{lll}
 \toprule
 Animal & Food & Size \\
 \midrule
 dog & meat & medium \\
 \cmidrule{1-2}
 horse & hay & large \\
 \cmidrule{1-1}
 \cmidrule{3-3}
 frog & flies & small \\
 \bottomrule
 \end{tabular}

\vspace{2em}

 \begin{tabular}{lll}
 \toprule
 Animal & Food & Size \\
 \midrule
 dog & meat & medium \\
 \cmidrule{1-2}
 horse & hay & large \\
 \cmidrule(r){1-1}
 \cmidrule(rl){2-2}
 \cmidrule(l){3-3}
 frog & flies & small \\
 \bottomrule
 \end{tabular}

\vspace{2em}

 \begin{tabular}{cp{9cm}}
 \toprule
 Animal & Description \\
 \midrule
 dog  & The dog is a member of the genus Canis, which forms part of the  wolf-like canids, and is the most widely abundant
 terrestrial carnivore. \\
  \addlinespace
 cat & The cat is a domestic species of small carnivorous mammal. It is the  only domesticated species in the family Felidae and is often referred  to as the domestic cat to distinguish it from the wild members of the family. \\
\bottomrule
 \end{tabular}


\subsection*{ 5.3 Merging cells}

 \begin{tabular}{lll}
 \toprule
 Animal & Food & Size \\
 \midrule
 dog & meat & medium \\
 horse & hay & large \\
 frog & flies & small \\
 fuath & \multicolumn{2}{c}{unknown} \\
 \bottomrule
 \end{tabular}

\vspace{2em}

 \begin{tabular}{lll}
 \toprule
 \multicolumn{1}{c}{Animal} & \multicolumn{1}{c}{Food} & \multicolumn{1}{c}{Size} \\
 \midrule
 dog & meat & medium \\
 horse & hay & large \\
 frog & flies & small \\
 fuath & \multicolumn{2}{c}{unknown} \\
 \bottomrule
 \end{tabular}

\vspace{2em}


 \begin{tabular}{lll}
 \toprule
 Group & Animal & Size \\
 \midrule
 herbivore 	& horse & large \\
 		& deer & medium \\
 		& rabbit & small \\
 \addlinespace
 carnivore 	& dog
 		& medium \\
		& cat & small \\
 		& lion & large \\
 \addlinespace
 omnivore 	& crow & small \\
		& bear & large \\
		& pig & medium \\
 \bottomrule
\end{tabular}

\vspace{2em}

\subsection*{ 5.4 The other preamble contents}

 \begin{tabular}{>{\itshape}l<{:} *{2}{l}}
 \toprule
 Animal & Food & Size \\
 \midrule
 dog & meat & medium \\
 horse & hay & large \\
 frog & flies & small \\
 \bottomrule
 \end{tabular}

\vspace{2em}


 \begin{tabular}{>{\itshape}l<{:} *{2}{l}}
 \toprule
 \multicolumn{1}{l}{Animal} & Food & Size \\
 \midrule
 dog & meat & medium \\
 horse & hay & large \\
 frog & flies & small \\
 \bottomrule
 \end{tabular}

\vspace{2em}


%\setlength\tabcolsep{1cm}
 \begin{tabular}{lll}
 Animal & Food & Size \\
 dog & meat & medium \\
 horse & hay & large \\
 frog & flies & small \\
 \end{tabular}

\vspace{2em}


 \begin{tabular}{l@{ : }l@{\hspace{2cm}}l}
 Animal & Food & Size \\
 dog & meat & medium \\
 horse & hay & large \\
 frog & flies & small \\
 \end{tabular}

\vspace{2em}

 \begin{tabular}{l!{:}ll}
 Animal & Food & Size \\
 dog
 & meat & medium \\
 horse & hay & large \\
 frog & flies & small \\
 \end{tabular}

\vspace{2em}



 \begin{tabular}{l|ll}
 Animal & Food & Size \\[2pt]
 dog & meat & medium \\
 horse & hay & large \\
 frog & flies & small \\
 \end{tabular}

\subsection*{ 5.5 Customizing booktabs rules}

 \begin{tabular}{@{} lll@{}} \toprule[2pt]
 Animal & Food & Size \\ \midrule[1pt]
 dog & meat & medium \\
 \cmidrule[0.5pt](r{1pt}l{1cm}){1-2}
 horse & hay & large \\
 frog & flies & small \\ \bottomrule[2pt]
 \end{tabular}

\subsection*{ 5.6 Numeric alignment in columns}

 \begin{tabular}{SS}
 \toprule
 {Values} & {More Values} \\
 \midrule
 1 	&  2.3456 \\
 1.2 	&  34.2345 \\
-2.3 	& 90.473 \\
 40 	& 5642.5 \\
 5.3 	& 1.2e3 \\
 0.2 	& 1e4 \\

 \bottomrule
 \end{tabular}

\subsection*{ 5.7 Specifying the total table width}



 \begin{center}
 \begin{tabular}{cc}
 \hline
 A & B\\
 C & D\\
 \hline
 \end{tabular}
 \end{center}
 \begin{center}
 \begin{tabular*}{.5\textwidth}{@{\extracolsep{\fill}}cc@{}}
 \hline
 A & B\\
 C & D\\
 \hline
 \end{tabular*}
 \end{center}
 \begin{center}
 \begin{tabular*}{\textwidth}{@{\extracolsep{\fill}}cc@{}}
 \hline
 A & B\\
 C & D\\
 \hline
 \end{tabular*}
 \end{center}

\vspace{2em}

 \begin{center}
 \begin{tabular}{lp{2cm}}
 \hline
 A & B B B B B B B B B B B B B B B B B B B B B B B B\\
 C & D D D D D D D\\
 \hline
 \end{tabular}
 \end{center}
 \begin{center}
 \begin{tabularx}{.5\textwidth}{lX}
 \hline
 A & B B B B B B B B B B B B B B B B B B B B B B B B\\
 C & D D D D D D D\\
 \hline
 \end{tabularx}
\end{center}
 \begin{center}
 \begin{tabularx}{\textwidth}{lX}
 \hline
 A & B B B B B B B B B B B B B B B B B B B B B B B B\\
 C & D D D D D D D\\
 \hline
 \end{tabularx}
 \end{center}


\vspace{2em}
\subsection*{ 5.8 Multi-page tables}

 \begin{longtable}{cc}
 \multicolumn{2}{c}{A Long Table}\\
 Left Side & Right Side\\
 \hline
 \endhead
 \hline
 \endfoot
 aa & bb\\
 Entry & b\\
 a & b\\
 a & b\\
 a & b\\
 a & b\\
 a & bbb\\
 a & b\\
 a & b\\
 a & b\\
 a & b\\
 a & b\\
 a & b\\
 a & b b b b b b\\
 a & b b b b b\\
 a & b b\\
 A Wider Entry & b\\
 \end{longtable}

\subsection*{ 5.9 Tablenotes}

 \begin{table}
 \begin{threeparttable}
 \caption{An Example}
 \begin{tabular}{ll}
 An entry & 42\tnote{1}\\
 Another entry & 24\tnote{2}\\
 \end{tabular}
 \begin{tablenotes}
 \item [1] the first note.
 \item [2] the second note.
 \end{tablenotes}
 \end{threeparttable}
 \end{table}

\subsection*{ 5.10 Typesetting in narrow columns}

 \begin{table}
 
 \begin{tabular}[t]{lp{3cm}}
 One & A long text set in a narrow paragraph, with some more example text.\\
 Two & A different long text set in a narrow paragraph, with some more hard to hyphenate words.

 \end{tabular}%
 \begin{tabular}[t]{l>{\raggedright\arraybackslash}p{3cm}}
 One & A long text set in a narrow paragraph, with some more example text.\\
 Two & A different long text set in a narrow paragraph, with some more hard to hyphenate words.

 \end{tabular}%
 \begin{tabular}[t]{l>{\RaggedRight}p{3cm}}
 One & A long text set in a narrow paragraph, with some more example text.\\
 Two & A different long text set in a narrow paragraph, with some more hard to hyphenate words.
 \end{tabular}
 \footnotesize
 \begin{tabular}[t]{lp{3cm}}
 One & A long text set in a narrow paragraph, with some more example text.\\
 Two & A different long text set in a narrow paragraph, with some more hard to hyphenate words.
 \end{tabular}
 \end{table}

\subsection*{ 5.12 Vertical tricks}

 \begin{tabular}{lcc}
 \toprule
 Test & \begin{tabular}{@{}c@{}}A\\a\end{tabular} & \begin{tabular}{@{}c@{}}B\\b\end{tabular} \\
 \midrule
 Content & is & here \\
 Content & is & here \\
 Content & is & here \\
 \bottomrule
 \end{tabular}

\vspace{2em}

 \begin{tabular}{lcc}
 \toprule
 Test & \begin{tabular}[b]{@{}c@{}}A\\a\end{tabular} &
 \begin{tabular}[t]{@{}c@{}}B\\b\end{tabular} \\
 \midrule
 Content & is & here \\
 Content & is & here \\
 Content & is & here \\
 \bottomrule
 \end{tabular}


\subsection*{ 5.13 Line spacing in tables}

 \renewcommand\arraystretch{1.5}
 \begin{center}
 \begin{tabular}{cc}
 \hline
 Square& $x^2$\\
 \hline
 Cube& $x^3$\\
 \hline
 \end{tabular}
 \end{center}

 \begin{center}


 %\setlength\extrarowheight{2pt}
 \begin{tabular}{cc}
 \hline
 Square& $x^2$\\
 \hline
 Cube& $x^3$\\
 \hline
 \end{tabular}
 \end{center}

\subsection*{ 5.14 Exercises}

\subsection*{ Use the simple table example to start experimenting with tables.}

 \begin{tabular}{lll}
 Animal & Food & Size \\
 dog & meat & medium \\
 horse & hay & large \\
 frog & flies & small \\
 \end{tabular}

\vspace{2em}

 \begin{tabular}{lll}
 Animal & Food  \\
 dog & meat  \\
 horse & hay  \\
 frog & flies \\
 \end{tabular}

\vspace{2em}

 \begin{tabular}{lll}
 Animal & Food & Size \\
 dog & meat & medium \\
 horse & hay & large \\
 \end{tabular}


\subsection*{ Try out different alignments using the l, c and r column types.}


 \begin{tabular}{lll}
 \toprule
 Animal & Food & Size \\
 \midrule
 dog & meat & medium \\
 \cmidrule{1-2}
 horse & hay & large \\
 \cmidrule{1-1}
 \cmidrule{3-3}
 frog & flies & small \\
 \bottomrule
 \end{tabular}

\vspace{2em}

 \begin{tabular}{ccc}
 \toprule
 Animal & Food & Size \\
 \midrule
 dog & meat & medium \\
 \cmidrule{1-2}
 horse & hay & large \\
 \cmidrule{1-1}
 \cmidrule{3-3}
 frog & flies & small \\
 \bottomrule
 \end{tabular}


 \begin{tabular}{rrr}
 \toprule
 Animal & Food & Size \\
 \midrule
 dog & meat & medium \\
 \cmidrule{1-2}
 horse & hay & large \\
 \cmidrule{1-1}
 \cmidrule{3-3}
 frog & flies & small \\
 \bottomrule
 \end{tabular}

\subsection*{What happens if you have too few items in a table row?}

 \begin{tabular}{lll}
A & B & C \\
1 & 2\\
 \end{tabular}

\subsection*{What happens if you have too many items in a table row?}

 \begin{tabular}{llll}
A & B & C \\
1 & 2 & 3 & 4
 \end{tabular}

\subsection*{Experiment with the multicolumn command to span across columns.}


 \begin{tabular}{lll}
 \toprule
 Animal & Food & Size \\
 \midrule
 dog & meat & medium \\
 horse & hay & large \\
 frog & flies & small \\
 fuath & \multicolumn{2}{c}{unknown} \\
 \bottomrule
 \end{tabular}

\vspace{2em}

 \begin{tabular}{lll}
 \toprule
 \multicolumn{3}{c}{Animal Information}\\
\hline
 Animal & Food & Size \\
 \midrule
 dog & meat & medium \\
 horse & hay & large \\
 frog & flies & small \\
 fuath & \multicolumn{2}{c}{unknown} \\
 \bottomrule
 \end{tabular}

\vspace{2em}

 \begin{tabular}{lll}
 \toprule
 Animal & Food & Size \\
 \midrule
 dog & meat & medium \\
 horse & hay & large \\
\multicolumn{3}{c}{unknown} \\
 fuath & \multicolumn{2}{c}{unknown} \\
 \bottomrule
 \end{tabular}

 \end{document}




